\section{CONCLUSION}
\label{sec:conclusion}

We benchmarked six trajectory-tracking controllers --- one per control paradigm, each
the strongest member of its own family in our data --- on a single nano-quadrotor, a
single reference family and a single command interface, across constant wind, gusts,
ground effect, a payload step and a measurement model of the optical positioning system
the vehicle actually carries, with per-cell seed statistics throughout.

No controller dominates, and the ranking is not a ranking of paradigms: everything that
beats a feedforward baseline under a sustained disturbance does so by estimating that
disturbance, and the paradigms differ in where the estimate enters rather than in
whether one is needed. The central result concerns the condition that matters most in
deployment. With realistic onboard sensing and wind,
an offset-free predictive controller, a reinforcement-learning policy trained with a
privileged critic, a classical observer stacked on a learned adaptive rate loop, and that
same observer driving the firmware attitude loop directly converge to within
\SI{6}{\milli\meter} of one another at approximately \SI{0.06}{\meter} --- each further
from its own clean-sensing performance than from the other three. When four unrelated
mechanisms saturate at the same value, the binding constraint lies elsewhere; we read
this regime as estimator-limited, which makes state estimation, not controller
architecture, the higher-leverage investment at this sensor quality.

One stack deserves separate comment. The reinforcement-learning policy is the weakest in
the pool under every clean-sensing condition we measured, and it reaches the leading
group only under degraded sensing --- the condition its training distribution was built
around. Read together with the randomization lesson below, this is the study's most
compact statement of what domain randomization buys and what it costs: not robustness
across conditions, but specialization to the conditions that were randomized over.

The failures proved more instructive than the ranking. A single observer bandwidth
cannot simultaneously filter noise and track gusts, and the obvious adaptive fix fails
because the scheduling statistic cannot separate an external disturbance from the
observer's own actuation-lag residual. Domain randomization does not make a memoryless
policy adaptive --- it makes it average over the randomization range --- and the cost is
measurable in clean-state agility. Layered adaptation composes when its layers address
disjoint error sources, rendering a \SI{23}{\percent} payload invisible, but the same
structure masks external forces from the observer above it. Two further lessons are
methodological, and we state them as corrections to our own earlier reporting: mean-only
tables concealed that two leading stacks fail only on unlucky sensor-bias draws while a
third never does, and two conclusions drawn from too few seeds --- a single-seed reading
that proved to be the minimum of ten draws, and a three-seed variance estimate that
overstated both mean and spread --- did not survive the seeds that should have been run
first.

The study has clear limits. The simulator omits aerodynamic drag, motor asymmetry and
battery sag, so absolute errors are optimistic and only relative comparisons are
claimed; a single airframe is used; and yaw is held fixed throughout. The hardware
campaign of Sec.~\ref{sec:method:hardware} addresses the first directly, testing three
pre-registered predictions in a programmable fan-array wind tunnel with motion-capture
ground truth, where the simulated disturbance classes are reproducible as physical set
points rather than merely approximated. Beyond it, the mechanisms identified here point
to specific next steps: recurrent or innovation-variance-aware policies so that a
randomized noise level becomes observable rather than averaged, and reference-aware
residual whitening so that an adaptive observer bandwidth can distinguish a disturbance
from its own lag.
